%============================================================
% CH1 · P2 · PROBLEM 2 — KEYS AND ROOMS (LeetCode 841)
%============================================================
\subsection{Problem 1.2.2: Keys and Rooms}
\problemlabel{1.2.2}{ssec:ch1-keys-rooms}

\problemcard{1.2.2}{Keys and Rooms}{LeetCode 841}{https://leetcode.com/problems/keys-and-rooms/}{\exmedium\ \lc{Medium}\ (the graph is hidden in the story)}{Model ``key to room $k$ found in room $r$'' as an edge $r \to k$; one search from room 0; count the reached rooms}

\subsubsection*{1. Problem statement}

There are $n$ rooms labelled $0..n-1$. All rooms are locked except room~$0$. Inside room $i$ lies a set of keys $\texttt{rooms}[i]$; a key with number $k$ opens room $k$, and you may carry every key you find. Return \texttt{true} if you can visit \emph{all} the rooms, \texttt{false} otherwise.

\textbf{Constraints:} $2 \le n \le 1000$, \; $0 \le |\texttt{rooms}[i]| \le 1000$, \; $1 \le \sum_i |\texttt{rooms}[i]| \le 3000$, \; $0 \le \texttt{rooms}[i][j] < n$, the keys inside one room are distinct.

\begin{table}[H]
\centering\small
\begin{tabular}{@{}lcl@{}}
\toprule
\textbf{rooms} & \textbf{Output} & \textbf{Why} \\
\midrule
$[[1],[2],[3],[\,]]$ & \texttt{true} & room 0 gives key 1, room 1 gives key 2, room 2 gives key 3 \\
$[[1,3],[3,0,1],[2],[0]]$ & \texttt{false} & the only key to room 2 lies inside room 2 \\
\bottomrule
\end{tabular}
\end{table}

\begin{diagrambox}
\begin{tikzpicture}[
    room/.style={rectangle, rounded corners=2pt, draw=deepblue, fill=theorybg, line width=0.9pt,
                 minimum width=1.5cm, minimum height=1.0cm, font=\small\bfseries, align=center},
    lockd/.style={room, draw=recursionred, fill=dupfill},
    key/.style={font=\scriptsize, text=accentorange}]
    \node[room] (r0) at (0,0) {room 0\\[-2pt]{\scriptsize\mdseries open}};
    \node[room] (r1) at (2.6,0) {room 1};
    \node[lockd] (r2) at (5.2,0) {room 2};
    \node[room] (r3) at (7.8,0) {room 3};
    \node[key] at (0,-0.85) {keys: 1, 3};
    \node[key] at (2.6,-0.85) {keys: 3, 0, 1};
    \node[key] at (6.3,-0.85) {keys: 2};
    \node[key] at (7.8,-0.85) {keys: 0};
    \draw[flow] (r0.north) to[bend left=25] node[lbl, above]{key 1} (r1.north);
    \draw[flow] (r0.north east) to[bend left=30] node[lbl, above]{key 3} (r3.north west);
    \draw[warningred, thick, -{Stealth[length=2.5mm]}] ([xshift=4mm]r2.south) to[out=-60, in=-120, looseness=4] node[lbl, below, text=warningred]{key 2} ([xshift=-4mm]r2.south);
\end{tikzpicture}
\figcap{Sample 2 as a building. Room 2 is red: its only key is locked inside it, so it can never be opened.}
\end{diagrambox}

\subsubsection*{2. The graph model}

The statement never says ``graph''. The modelling step is to notice that \emph{finding key $k$ in room $r$} means \emph{``once I am in $r$, I can go to $k$''} --- an edge $r \to k$.

\begin{modelbox}[Graph Model --- Keys and Rooms]
\renewcommand{\arraystretch}{1.35}
\begin{tabularx}{\linewidth}{@{}>{\bfseries}l X@{}}
Vertices & the $n$ rooms \\
Edges & directed: $r \to k$ for every key $k \in \texttt{rooms}[r]$. The input \emph{is} the adjacency list --- nothing to build. Self-loops are possible (room 1 of Sample~2 contains key 1) and harmless \\
Weights & none \\
Question & is $R(0) = \{0, 1, \dots, n-1\}$? i.e.\ does one search from room 0 mark all $n$ vertices? \\
Size & $|V| \le 1000$, $|E| = \sum|\texttt{rooms}[i]| \le 3000$: tiny. Even the $\bigO{|V|\cdot(|V|+|E|)}$ brute force below fits; the point is to see the graph \\
\end{tabularx}
\end{modelbox}

\begin{diagrambox}
\begin{tikzpicture}
    \begin{scope}
        \node[gvsrc] (a0) at (0,0) {0}; \node[gv] (a1) at (1.4,0) {1};
        \node[gv] (a2) at (2.8,0) {2}; \node[gv] (a3) at (4.2,0) {3};
        \draw[gde] (a0) -- (a1); \draw[gde] (a1) -- (a2); \draw[gde] (a2) -- (a3);
        \node[note] at (2.1,-0.9) {Sample 1: a chain $0\to1\to2\to3$, \texttt{true}};
    \end{scope}
    \begin{scope}[xshift=6.6cm]
        \node[gvsrc] (b0) at (0,0.9) {0}; \node[gv] (b1) at (2.0,0.9) {1};
        \node[gv] (b3) at (1.0,-0.7) {3}; \node[gvdst] (b2) at (3.8,0) {2};
        \draw[gde] (b0) to[bend left=15] (b1); \draw[gde] (b1) to[bend left=15] (b0);
        \draw[gde] (b0) to[bend left=15] (b3); \draw[gde] (b3) to[bend left=15] (b0);
        \draw[gde] (b1) -- (b3);
        \draw[gde] (b1) to[out=60, in=120, looseness=6] (b1);
        \draw[gde] (b2) to[out=60, in=120, looseness=6] (b2);
        \draw[gcut] (2.9,1.7) -- (2.9,-1.1);
        \node[note] at (1.9,-1.45) {Sample 2: no edge enters 2 from outside, \texttt{false}};
    \end{scope}
\end{tikzpicture}
\figcap{Left: the rooms of Sample 1. Right: Sample 2. Room 2's only incoming edge is its own self-loop; the red line separates $\{0,1,3\}$, which has no edge leaving it.}
\end{diagrambox}

\subsubsection*{3. Key observations}

\begin{itemize}[leftmargin=*]
    \item \textbf{Keys are never used up and doors never close.} Holding more keys never hurts, so the \emph{order} in which we open rooms does not matter. The set of rooms we can eventually enter is fixed --- it is $R(0)$. (Contrast: if a key disappeared when used, the order would matter and this would be a much harder search over states.)
    \item \textbf{``Visit all'' $=$ ``everything reachable''.} Count the rooms marked by the search and compare with $n$.
    \item \textbf{Direction matters.} Key $k$ in room $r$ helps only after reaching $r$: edge $r \to k$, never $k \to r$. Room 3 of Sample~2 holds key 0, which is useless --- room 0 was open anyway.
    \item \textbf{A local check is not enough.} ``Every room except 0 has a key somewhere'' holds in Sample~2 (key 2 is in room 2) and still the answer is \texttt{false}: the key has to lie in a room you can \emph{reach}.
\end{itemize}

\subsubsection*{4. Approaches}

\begin{bruteforcebox}
\textbf{Idea: simulate a visitor who repeatedly walks along the corridor.} Each pass walks past rooms $0, 1, \dots, n-1$ and enters every room for which a key is in hand and that has not been entered yet, picking up its keys. Stop when a whole pass enters nothing.

\begin{lstlisting}[style=cpp]
#include <bits/stdc++.h>
using namespace std;

class Solution {
public:
    bool canVisitAllRooms(vector<vector<int>>& rooms) {
        int n = rooms.size();
        vector<char> haveKey(n, 0), entered(n, 0);
        haveKey[0] = 1;                                // room 0 is unlocked
        bool changed = true;
        while (changed) {                              // one walk along the corridor
            changed = false;
            for (int r = 0; r < n; r++)
                if (haveKey[r] && !entered[r]) {
                    entered[r] = 1; changed = true;
                    for (int k : rooms[r]) haveKey[k] = 1;
                }
        }
        for (int r = 0; r < n; r++) if (!entered[r]) return false;
        return true;
    }
};
\end{lstlisting}

Correct by the same closure argument as the pattern, but a pass can enter a single room. Worst case: room 0 holds key $n{-}1$, and room $i \ge 2$ holds key $i{-}1$. With $n = 6$ the program printed:
\begin{center}\small
\begin{tabular}{@{}lcccccc@{}}
\toprule
pass & 1 & 2 & 3 & 4 & 5 & 6 \\
rooms entered & 0, 5 & 4 & 3 & 2 & 1 & --- \\
\bottomrule
\end{tabular}
\end{center}
\begin{center}
\begin{tikzpicture}[x=1.25cm]
    \foreach \i/\p in {0/1, 1/5, 2/4, 3/3, 4/2, 5/1}{
        \ifnum\i=0 \node[gvsrc] (q\i) at (\i,0) {\i}; \else \node[gvdone] (q\i) at (\i,0) {\i}; \fi
        \node[gdist] at (\i,-0.6) {pass \p};
    }
    \draw[gde] (q0) to[bend left=30] (q5);
    \foreach \i [evaluate=\i as \j using int(\i-1)] in {5,4,3,2} \draw[gde] (q\i) to[bend left=30] (q\j);
    \draw[->, thick, darkgray] (-0.4,-1.05) -- node[note, below] {direction of every corridor walk} (5.4,-1.05);
\end{tikzpicture}
\end{center}
The key chain $0 \to 5 \to 4 \to 3 \to 2 \to 1$ runs \emph{against} the corridor: each key is found just after the walk has passed its door, so for $n = 1000$ the program needed $1000$ passes.

\begin{complexitybox}
\textbf{Time:} $\bigO{n \cdot (n + \sum|\texttt{rooms}[i]|)}$ --- about $4 \cdot 10^6$ here, so it is even accepted, but it is quadratic: with $n = 10^5$ rooms it would be hopeless. \textbf{Space:} $\bigO{n}$. \textbf{Verdict:} \ac{} only thanks to the small limits.
\end{complexitybox}
\end{bruteforcebox}

\begin{algorithmbox}[Approach 2: DFS with a stack of unused keys]
\textbf{Idea.} Instead of walking past every door again and again, keep a \textbf{pile of keys for rooms we have not entered yet}. Take any key from the pile, enter that room, and add its new keys to the pile. This is the pattern's search: the pile is the to-do stack, ``room opened'' is \texttt{seen}.

\begin{lstlisting}[style=cpp]
#include <bits/stdc++.h>
using namespace std;

class Solution {
public:
    bool canVisitAllRooms(vector<vector<int>>& rooms) {
        int n = rooms.size();
        vector<char> opened(n, 0);
        vector<int> pile(1, 0);                       // keys for rooms not yet entered
        opened[0] = 1;
        int count = 1;                                // rooms opened so far
        while (!pile.empty()) {
            int r = pile.back(); pile.pop_back();     // enter room r
            for (int k : rooms[r])                    // edge r -> k
                if (!opened[k]) { opened[k] = 1; count++; pile.push_back(k); }
        }
        return count == n;
    }
};
\end{lstlisting}

\begin{complexitybox}
\textbf{Time:} $\bigO{n + \sum|\texttt{rooms}[i]|}$ --- every room is entered once, every key is looked at once. \textbf{Space:} $\bigO{n}$ (the input already is the adjacency list). \textbf{Verdict:} \ac{}.
\end{complexitybox}
\end{algorithmbox}

\begin{invariantbox}[Why it is correct]
\textbf{Invariant.} At any moment, \texttt{opened} $=$ the rooms entered so far $\cup$ the rooms whose key is on the pile, and every opened room can really be entered.

\emph{Soundness.} Room 0 can be entered. A room $k$ is marked only when its key is found in a room $r$ we actually entered, so $0 \leadsto r \to k$.

\emph{Completeness.} When the pile is empty, every opened room has been entered and all its keys were examined, so no key for an unopened room lies in an opened room --- the opened set has no edge leaving it. Any room reachable from 0 is therefore opened (the pattern's closure argument). Since \texttt{count} is the number of opened rooms, \texttt{count == n} holds exactly when $R(0)$ contains every room. \hfill$\blacksquare$
\end{invariantbox}

\begin{alternativebox}[Approach 3: BFS or recursive DFS --- same set, different order]
Replacing the pile by a queue visits the rooms in order of ``how many doors away'' they are; on Sample~2 the BFS order is $0, 1, 3$ and the DFS order is $0, 3, 1$ --- the same \emph{set} $\{0,1,3\}$, so the same answer.

\begin{center}
\begin{tikzpicture}
    \foreach \x/\name/\o in {0/BFS (queue)/{0, 1, 3}, 5.5/DFS (stack)/{0, 3, 1}}{
        \begin{scope}[xshift=\x cm]
            \node[gvsrc] (h0) at (0,0.9) {0}; \node[gvdone] (h1) at (2.0,0.9) {1};
            \node[gvdone] (h3) at (1.0,-0.7) {3}; \node[gvdead] (h2) at (3.2,0) {2};
            \draw[g1p2tree] (h0) -- (h1); \draw[g1p2tree] (h0) -- (h3);
            \draw[gfaded, -{Stealth[length=2.2mm]}] (h1) -- (h3);
            \node[font=\small\bfseries] at (1.6,1.7) {\name: order \o};
        \end{scope}
    }
\end{tikzpicture}
\end{center}
Here even the search trees coincide (both consist of the two edges out of room 0); only the visiting order differs. The recursive version is the shortest to write and safe here, because the recursion depth is at most $n \le 1000$:

\begin{lstlisting}[style=cpp]
#include <bits/stdc++.h>
using namespace std;

class Solution {
    vector<char> opened;
    int count = 0;
    void enter(int r, vector<vector<int>>& rooms) {
        opened[r] = 1; count++;
        for (int k : rooms[r])
            if (!opened[k]) enter(k, rooms);
    }
public:
    bool canVisitAllRooms(vector<vector<int>>& rooms) {
        opened.assign(rooms.size(), 0);
        count = 0;
        enter(0, rooms);
        return count == (int)rooms.size();
    }
};
\end{lstlisting}
\begin{complexitybox}
\textbf{Time:} $\bigO{n + \sum|\texttt{rooms}[i]|}$. \textbf{Space:} $\bigO{n}$ plus a recursion depth $\le n$. \textbf{Verdict:} \ac{} --- but with $n = 10^6$ only the iterative versions would be safe.
\end{complexitybox}
\end{alternativebox}

\begin{tricksbox}[Approach 4: Implementation Refinements]
\begin{itemize}[leftmargin=*]
    \item \textbf{No graph building.} When the input is already ``for each vertex, its out-neighbours'', use it directly as \texttt{adj}. Recognising that saves both code and memory.
    \item \textbf{Early success.} One may stop as soon as \texttt{count == n}; it does not change the complexity but helps when most keys are in the first rooms.
    \item \textbf{Report the locked rooms.} After the search, every $r$ with \texttt{opened[r] == 0} is a room that can never be entered --- useful for debugging and for variants that ask ``which rooms?''.
\end{itemize}
\end{tricksbox}

\subsubsection*{5. Trace}

\begin{tracebox}[{Sample 2: rooms $= [[1,3],[3,0,1],[2],[0]]$}]
\begin{center}
\begin{tikzpicture}[scale=0.72, transform shape]
    \foreach \k/\t/\a/\b/\c/\d in {
        0/{start: pile $[0]$}/gvsrc/gv/gv/gv,
        1/{enter 0: pile $[1,3]$}/gvdone/gvseen/gv/gvseen,
        2/{enter 3: pile $[1]$}/gvdone/gvseen/gv/gvdone,
        3/{enter 1: pile empty}/gvdone/gvdone/gvdead/gvdone}{
        \begin{scope}[xshift=\k*5.4cm]
            \node[\a] (z0) at (0,0.9) {0}; \node[\b] (z1) at (2.0,0.9) {1};
            \node[\d] (z3) at (1.0,-0.7) {3}; \node[\c] (z2) at (3.3,0) {2};
            \draw[gde] (z1) to[bend left=15] (z0);
            \draw[gde] (z3) to[bend left=15] (z0);
            \draw[gde] (z1) -- (z3);
            \draw[gde] (z1) to[out=60, in=120, looseness=6] (z1);
            \draw[gde] (z2) to[out=60, in=120, looseness=6] (z2);
            \ifnum\k>0 \draw[g1p2tree] (z0) to[bend left=15] (z1); \draw[g1p2tree] (z0) to[bend left=15] (z3);
            \else \draw[gde] (z0) to[bend left=15] (z1); \draw[gde] (z0) to[bend left=15] (z3); \fi
            \node[font=\large\bfseries] at (1.6,2.2) {\t};
        \end{scope}
    }
\end{tikzpicture}
\end{center}
\begin{center}\small
\begin{tabular}{@{}cllcl@{}}
\toprule
\textbf{step} & \textbf{enter} & \textbf{keys found $\to$ new rooms} & \textbf{count} & \textbf{pile after (top on the right)} \\
\midrule
1 & 0 & $1, 3 \to$ 1, 3 & 3 & 1 3 \\
2 & 3 & $0 \to$ --- & 3 & 1 \\
3 & 1 & $3, 0, 1 \to$ --- & 3 & (empty) \\
\bottomrule
\end{tabular}
\end{center}
$\texttt{count} = 3 \ne 4 \Rightarrow$ \texttt{false}. Room 1's self-key (1) and room 3's key 0 are simply skipped by the \texttt{opened} test. Room 2 is never entered, so its self-loop is never even looked at.
\end{tracebox}

\begin{tracebox}[{Sample 1: rooms $= [[1],[2],[3],[\,]]$}]
\begin{center}
\begin{tikzpicture}[scale=0.8, transform shape]
    \foreach \k/\t/\s in {0/{enter 0, pile $[1]$}/1, 1/{enter 1, pile $[2]$}/2, 2/{enter 2, pile $[3]$}/3, 3/{enter 3, pile empty}/4}{
        \begin{scope}[xshift=\k*4.3cm]
            \foreach \i in {0,...,3}{
                \pgfmathtruncatemacro{\cur}{\s-1}
                \ifnum\i<\cur \node[gvdone] (w\i) at (\i*0.95,0) {\i};
                \else\ifnum\i=\cur \node[gvcur] (w\i) at (\i*0.95,0) {\i};
                \else\ifnum\i=\s \node[gvseen] (w\i) at (\i*0.95,0) {\i};
                \else \node[gv] (w\i) at (\i*0.95,0) {\i};\fi\fi\fi
            }
            \foreach \i [evaluate=\i as \j using int(\i+1)] in {0,...,2}{
                \ifnum\i<\s \draw[g1p2tree] (w\i) -- (w\j); \else \draw[gde] (w\i) -- (w\j); \fi}
            \node[font=\small\bfseries] at (1.45,0.75) {\t};
        \end{scope}
    }
\end{tikzpicture}
\end{center}
Every room yields exactly one new key; the pile never holds more than one key. After room 3, $\texttt{count} = 4 = n \Rightarrow$ \texttt{true}. The search tree is the whole chain.
\end{tracebox}

\subsubsection*{6. Test suite}

\begin{testbox}
\centering\small
\begin{tabular}{@{}p{6.4cm}cp{5.2cm}@{}}
\toprule
\textbf{rooms} & \textbf{Expected} & \textbf{Why} \\
\midrule
$[[1],[\,]]$ & \texttt{true} & smallest \texttt{true} \\
$[[\,],[0]]$ & \texttt{false} & key exists, but in the wrong room (direction) \\
$[[0],[0]]$ & \texttt{false} & self-key in room 0; room 1 unreachable \\
chain $i \to i+1$, $n = 1000$ & \texttt{true} & longest possible search path \\
room 0 holds keys $1..999$ & \texttt{true} & one room with 999 out-edges \\
chain $0 \to \dots \to 998 \to 0$, room 999 holds key 0 & \texttt{false} & the last room's only edge goes \emph{out} of it \\
\bottomrule
\end{tabular}
\end{testbox}
The expected values come from running Approach~2; in addition, Approaches 1--3 agreed with a Floyd--Warshall closure on 30\,000 random instances with $2 \le n \le 8$.

\subsubsection*{7. Common mistakes}

\begin{warningbox}[Mistake 1 --- edges in the wrong direction]
Building \texttt{adj[k].push\_back(r)} (``room $k$ is opened \emph{by} room $r$'') answers the reverse question ``which rooms lead to room 0?''. On the last test above it would mark every room and wrongly return \texttt{true}.
\end{warningbox}

\begin{warningbox}[Mistake 2 --- counting keys instead of rooms]
\begin{lstlisting}[style=snippet]
for (int k : rooms[r]) { count++; ... }         // WRONG: every key counted, even old ones
for (int k : rooms[r]) if (!opened[k]) { opened[k] = 1; count++; ... }   // RIGHT
\end{lstlisting}
On $[[1],[0]]$ the wrong count is $1 + 1 + 1 = 3 \ne 2$, so it answers \texttt{false} although both rooms open; on Sample~2 it reaches 7 (keys 3, 0 and 1 are counted twice). Only \emph{newly opened} rooms may be counted.
\end{warningbox}

\begin{warningbox}[Mistake 3 --- forgetting that room 0 is already open]
Starting with \texttt{count = 0} and \texttt{opened[0] = 0} makes room 0 count only if some room contains key 0. Room 0 must be marked and counted before the loop.
\end{warningbox}

\begin{warningbox}[Mistake 4 --- ``every room has a key somewhere'']
Checking that each room $1..n{-}1$ appears in some key list is a Pattern~1 (in-degree $>0$) test. Sample~2 passes it and is still \texttt{false}: in-degree from \emph{reachable} rooms is what matters, and that needs a search.
\end{warningbox}

\subsubsection*{8. Takeaways}

\begin{itemize}[leftmargin=*]
    \item ``Finding $X$ in $A$ lets you use $B$'' is an edge $A \to B$; ``can we get all of them?'' is ``$|R(\text{start})| = n$?''.
    \item Monotone processes (keys never vanish, infections never heal) are pure reachability: order does not matter.
    \item Sometimes the input already is the adjacency list --- recognise it and skip the building step.
    \item Count \emph{newly marked} vertices, not edges.
\end{itemize}

\begin{variantsbox}
\variant{Fewest extra unlocked rooms}{\varin{https://leetcode.com/problems/minimum-number-of-vertices-to-reach-all-nodes/}}
  {choose as few starting rooms as possible so that every room becomes reachable.}
  {same graph; we want the smallest set $S$ with $\bigcup_{s\in S} R(s) = V$.}
  {if the graph is a DAG: exactly the vertices of in-degree 0 (E2.2); in general one per ``source'' strongly connected component (Chapter~19).}
\variant{How many rooms from every start?}{\varin{https://atcoder.jp/contests/abc204/tasks/abc204_c}}
  {for every room as the start, count the rooms reachable from it (summed over all starts).}
  {the same directed graph with $|V| \le 2000$, $|E| \le 2000$.}
  {one search per start vertex, $\bigO{|V|(|V|+|E|)}$ (E2.3).}
\variant{Several rooms open initially}{\varnone}
  {rooms in a set $S$ are unlocked at the start, not only room 0.}
  {same graph, many sources.}
  {multi-source search: push every room of $S$ before the loop; still $\VE$.}
\variant{Which rooms stay locked?}{\varnone}
  {output the rooms that can never be entered.}
  {same graph.}
  {the same search; output every $r$ with \texttt{opened[r] == 0} (the complement of $R(0)$).}
\variant{Keys break after one use}{\varnone}
  {each key opens one door and is then lost; rooms may need to be re-entered.}
  {the state is no longer ``which room'' but ``which keys are in hand'' --- an implicit state graph.}
  {search over states (Pattern~4); the order of moves now matters, and the state space can be exponential.}
\end{variantsbox}
